%% ==========================================================================
%% SECTION 1: INTRODUCTION (RESTRUCTURED)
%% ==========================================================================
\section{Introduction}\label{sec:intro}

Systems that incorporate language model components face a verification problem.
The components generate text by sampling from learned probability distributions,
and some fraction of that text is fabricated: fluent, confident, and wrong. We
prefer the term \emph{fabrication} throughout this paper to emphasize that
fluent confabulation is the system's default behavior under coherence
optimization, not a malfunction to be patched. A language model optimized for
coherence produces coherent output. When the training data provides grounding,
that output tends to track truth. When it does not, the same optimization
produces fluent, confident fabrication. Xu et
al.~\cite{xu2024hallucination} prove a related result: hallucination is
inevitable for LLMs over the class of computable functions. Our result is
orthogonal. They prove a computational limit on what the model can produce;
we prove an observational limit on what a supervisor can verify. The model may
know it is fabricating. The supervisor cannot tell.

\fakepara{Verification costs money.}
Checking a citation against a database costs API calls and latency. Having a
domain expert review a medical summary costs time and salary. Running a second
model as a judge costs compute. Every system that deploys a language model
component makes, explicitly or implicitly, a budget allocation decision: how
much of the output to verify, which outputs to prioritize, and what signals to
use for triage. But the standard text-only interface discards the very signals
that would make verification efficient. A supervisor receiving only text must
spend its verification budget reconstructing, from the output alone, what the
model's own computation already knew. The budget is consumed by reconstruction
rather than concentrated on the cases where verification is most needed.

\fakepara{Self-report is worse than useless.}
Across four model architectures (OLMo-3 7B, Llama-3.1 8B, Qwen3 4B,
Mistral 7B), every model reports \emph{higher} confidence on fabrications than
on known facts. Self-report AUC ranges 0.28--0.36---below random, because the
relationship is systematically inverted. The models carry the epistemic state
needed for verification. The interface throws it away. Stacking supervisors
does not escape the constraint---later judges see no finer information than
earlier ones.

\fakepara{Testimony versus telemetry.}
Language models carry internal signals---entropy traces, attention geometry,
token-level log-probabilities---that distinguish grounded generation from
fabrication. These signals are byproducts of the computation itself:
\emph{telemetry} (measurements of the computation that produced the output)
that the model cannot separately control under standard training. But the
standard text-only interface exports only the model's chosen output---its
\emph{testimony} (what the model says). The testimony may lie; under current
training objectives, the telemetry is harder to fake.
Code plagiarism detection systems like
MOSS~\cite{schleimer2003winnowing, mason2019collaboration} have navigated this
same structural distinction for decades. MOSS works because gaming the
similarity metric requires writing genuinely different code---restructuring
code to defeat the detector changes the code itself. Behavioral signals
(formatting, variable naming, comment style) can be trivially altered without
changing the underlying logic, but computational signals (algorithmic structure,
data flow) cannot be independently controlled. The same pattern governs
epistemic verification of language model outputs: behavioral signals (response
length, hedging frequency) can be adapted by training without altering what the
model knows, while computational signals (per-token entropy, attention
geometry) cannot be flattened without changing the generation process itself.

\fakepara{A tensor interface.}
We construct a minimal interface extension that exports computational
telemetry alongside text, with no architectural changes to current
transformers. The interface exports entropy traces and attention
summaries---telemetry that the model cannot independently control. Under
bounded verification, tensor-guided judges at 10\% budget outperform
text-guided judges at 30\% budget (82.1\% vs.\ 78.5\% accuracy), not because
the tensor signal is strong in isolation, but because it enables efficient
triage: concentrating verification on the cases most likely to contain
fabrications.

\fakepara{Three results.}
This paper establishes three results that together characterize the cost
surface of epistemic verification.

First, we prove an \emph{impossibility}: epistemic honesty is unverifiable for
systems operating under text-only observation with bounded supervision
(\autoref{sec:formal}). The impossibility decomposes into two theorems. A
predictor-centric policy conditioning only on the query cannot satisfy
epistemic honesty across ambiguous world states
(\autoref{thm:representational}). Even with retrieval augmentation, a bounded
supervisor cannot provide the training signal needed to learn honest behavior
because plausible fabrications are observationally indistinguishable from
truthful responses (\autoref{thm:learnability}). This impossibility functions
as a budget constraint: below this investment line, text-channel verification
cannot help.

Second, we construct a \emph{tensor interface} that escapes the impossibility
class by exporting computational signals alongside text
(\autoref{sec:escape}). The interface requires no architectural changes to
current transformers.

Third, we provide an \emph{honest decomposition} of where the discrimination
comes from (\autoref{sec:eval}). The text channel is stronger than
expected: response length alone achieves AUC 0.88--0.96 across models---a
free signal requiring no tensor interface. But it is trivially gameable: a
model trained to fabricate briefly defeats length-based detection. After
controlling for length, the tensor residual adds AUC 0.60--0.70: a
computational signal that cannot be independently gamed without adversarial
optimization targeting the signal itself. The gap between the text-channel
ceiling and the tensor-channel floor is the impossibility theorem made
empirical---it measures the boundary between gameable and robust verification.

\fakepara{A cost map, not a detector.}
The contribution of this paper is not a high AUC number. It is the map.
Preliminary observations suggest verification cost varies by content format:
natural text, documented code, compact code, and mathematical proofs exhibit
different scaffolding ratios, predicting different optimal judge strategies.
Systems engineers building with language model components need to know what
each level of verification investment buys, and where the investment
transitions from text-channel (cheap, gameable) to tensor-channel (more
expensive, robust). This paper provides that cost surface.

This paper makes the following contributions:

\begin{itemize}
\item \textbf{An impossibility result as budget constraint.}
  \autoref{sec:formal} proves that
  epistemic honesty is unverifiable for predictive token generation
  systems operating under text-only observation with bounded
  supervision. The impossibility is structural: it holds regardless
  of model scale, training procedure, or prompt engineering, and
  cannot be escaped by stacking supervisors. It establishes the
  floor below which text-channel verification investment cannot improve
  assurance.

\item \textbf{A tensor interface as the next investment tier.}
  \autoref{sec:escape} demonstrates a minimal interface extension that exports
  epistemic telemetry alongside text with no architectural changes to
  current transformers. The tensor interface increases the observation
  budget available to external verifiers.

\item \textbf{An empirical cost surface.}
  \autoref{sec:eval} decomposes verification effectiveness into
  text-channel and tensor-channel components, characterizes the
  cost-robustness tradeoff, and identifies the format-dependent
  variation that determines judge strategy per domain.
\end{itemize}
